% Source: https://artofproblemsolving.com/community/c3h2745929_solving_an_equation_in_different_ways
% Stable identity: aops-topic-2745929-q01
Isolating $k$: $k(\sqrt5-1)=4$, so
\[k=\dfrac{4}{\sqrt5-1}=\dfrac{4(\sqrt5+1)}{4}=1+\sqrt5.\]
This method uses only reversible steps, so it gives exactly the solutions.

Squaring is not reversible: $(k+4)^2=5k^2$ also holds when $k+4=-k\sqrt5$. It produces $k=1\pm\sqrt5$, and one must check them. For $k=1-\sqrt5$: $k+4=5-\sqrt5>0$ but $k\sqrt5=\sqrt5-5<0$, so it is extraneous; it solves $k+4=-k\sqrt5$ instead. The only solution is $\boxed{k=1+\sqrt5}$.
